%% Example CV - Pablo Santana Susilla
%% Comprehensive overview of all competencies, experience, and achievements.
%% Use as a master reference when tailoring role-specific CVs.
%%
%% SETUP: Run /setup to populate this with your actual information.
%% Compile with: lualatex main_example.tex
%% (pdflatex often fails on modern MiKTeX with fontawesome5 font-expansion errors.)

\documentclass[11pt,a4paper,sans]{moderncv}
\moderncvstyle{banking}
\moderncvcolor{blue}

% Force the name and section headings to render in moderncv blue (color1).
% Default banking leaves them black: moderncvstylebanking.sty's \colorlet
% copies (not aliases) the pre-scheme accent colour, so the name colours are
% frozen before \moderncvcolor runs. Re-let them after. \namefont is the hook
% every name-style macro routes through, so this also works on moderncv 2.3.1
% (Debian/Ubuntu apt), which has no \firstnamestyle/\lastnamestyle at all.
\renewcommand*{\namefont}{\fontsize{34}{36}\bfseries\upshape}
\colorlet{firstnamecolor}{color1}
\colorlet{lastnamecolor}{color1}
\colorlet{namecolor}{color1}
\renewcommand*{\sectionstyle}[1]{{\sectionfont\color{color1}#1}}

\usepackage[utf8]{inputenc}
% pdflatex fallback only (the documented engine is lualatex, which skips this
% branch). Without T1 font encoding pdflatex builds accented letters with
% \accent, and the PDF text layer stores them decomposed - `e` + U+0300 rather
% than U+00E8 - so an ATS keyword match on "Genève" fails while the page looks
% right. moderncv 2.5 loads T1 itself under pdflatex; 2.3.1 (Debian/Ubuntu apt)
% does not. \ifpdftex comes from iftex, which every moderncv version loads.
\ifpdftex\usepackage[T1]{fontenc}\fi
% moderncv loads hyperref itself in an \AtEndPreamble hook, so \hypersetup
% must go in an \AtEndPreamble of our own: on moderncv < 2.4 a top-level
% \usepackage{hyperref} clashes with the class's own
% \RequirePackage[unicode]{hyperref}. From 2.4.0 the class passes its options
% through \PassOptionsToPackage instead, which is what removes that clash.
\AtEndPreamble{\hypersetup{
    colorlinks=true,
    linkcolor=blue,
    filecolor=magenta,
    urlcolor=blue,
    pdftitle={Pablo Santana Susilla - CV},
    % Keep pdfpagemode=UseNone: this block runs after moderncv's own
    % \AtEndPreamble (moderncv.cls sets pdfpagemode there), so a FullScreen
    % value here would win and open every CV in fullscreen presentation mode.
    pdfpagemode=UseNone,
}}
\usepackage[scale=0.80]{geometry}
\usepackage{import}

% personal data
\name{Pablo}{Santana Susilla}
% If you have no address to list, DELETE this whole line. \address{}{}{} fails
% with "There's no line here to end" on every moderncv version.
\address{Las Palmas de Gran Canaria, Canarias, España}{}{}
\email{pasantanasu@hotmail.com}
\extrainfo{\href{https://www.linkedin.com/in/pablo-santana-susilla-092248314}{LinkedIn}}

\begin{document}

\makecvtitle

% ============================================================
%     PROFILE STATEMENT
% ============================================================

\vspace{6pt}
\small{Desarrollador Odoo e ingeniero informático que construye módulos a medida desde cero y los mantiene en producción. Desarrollé la primera versión del módulo de Expedientes de subvenciones, núcleo del ERP Odoo sobre el que opera hoy toda Orbalia (25 personas), y diseñé la migración completa desde ActiveCampaign (más de 7.000 expedientes, 10.000 contactos y 1.200 empresas). Recojo requisitos con negocio, diseño el modelo de datos, desarrollo, despliego en AWS y mantengo. Antes, consultor tecnológico para pymes en el programa Kit Consulting.}

% ============================================================
%     CORE COMPETENCIES
% ============================================================

\section{Competencias clave}
\vspace{1pt}
\begin{itemize}

\item \textbf{Odoo}: desarrollo de módulos a medida (ORM, herencia, campos computados), vistas XML, informes QWeb, reglas de seguridad y permisos.

\item \textbf{Desarrollo e integración}: Python, JavaScript, PostgreSQL, API REST, XML-RPC, migración de datos.

\item \textbf{Infraestructura}: administración de producción en AWS (despliegues, actualizaciones, disponibilidad), Linux, Git.

\item \textbf{Dominio}: implantación de ERP, gestión de subvenciones, cumplimiento normativo (Ley de Transparencia de Canarias), automatización contable.

\item \textbf{Consultoría}: recogida de requisitos, diseño de modelos de datos, selección e implantación de soluciones digitales para pymes.

\end{itemize}

% ============================================================
%     PROFESSIONAL EXPERIENCE
% ============================================================

\section{Experiencia profesional}
\vspace{3pt}
\begin{itemize}

\item{\cventry{2025--Actualidad}{Consultor IT -- Desarrollador Odoo}{Orbalia}{Las Palmas de Gran Canaria}{}{\vspace{1pt}
\begin{itemize}
    \item {Desarrollé la primera versión del módulo de Expedientes de subvenciones (Odoo, Python, XML), núcleo del ERP que sustituyó a ActiveCampaign y sobre el que opera toda la empresa (25 personas); más de 230 commits.}
    \item {Diseñé y ejecuté la migración completa desde ActiveCampaign con scripts en Python (API REST a XML-RPC de Odoo): más de 7.000 expedientes, 10.000 contactos y 1.200 empresas sin perder relaciones.}
    \item {Diseñé y desarrollé en solitario el módulo de Proyectos IT, la detección automática de clientes sujetos a la Ley de Transparencia de Canarias y la creación automática de la cuenta contable de cada cliente.}
    \item {Definí con el equipo los modelos transversales (Contactos, Clientes, Grupos empresariales) y codesarrollé los módulos de Financiación y Deducciones fiscales.}
    \item {Administro la infraestructura en AWS (despliegues, actualizaciones, disponibilidad) como una de las dos personas con acceso al entorno.}
\end{itemize}}}

\vspace{3pt}

\item{\cventry{2025}{Consultor en formación (Kit Consulting)}{Orbalia}{Las Palmas de Gran Canaria}{}{\vspace{1pt}
\begin{itemize}
    \item {Asesoré a pymes en el diagnóstico de sus necesidades y en la selección e implantación de soluciones digitales.}
    \item {Segunda persona incorporada al departamento de IT, creado en septiembre de 2024.}
\end{itemize}}}

\vspace{3pt}

\item{\cventry{2024}{Desarrollador web}{Tca-tik}{Las Palmas de Gran Canaria}{}{\vspace{1pt}
\begin{itemize}
    \item {Desarrollo web con Laravel y JavaScript.}
\end{itemize}}}

\end{itemize}

% ============================================================
%     EDUCATION
% ============================================================

\section{Formación}
\vspace{1pt}
\begin{itemize}

\item{\cventry{2018--2024}{Grado en Ingeniería Informática}{Universidad de Las Palmas de Gran Canaria}{Las Palmas de Gran Canaria}{}{\vspace{1pt}
}}

\vspace{3pt}

\item{\cventry{2021--2022}{Estudios de Ingeniería Informática (BTech)}{Vilnius University}{Vilna, Lituania}{}{\vspace{1pt}
}}

\vspace{3pt}

\item{\cventry{Certificaciones}{CCNAv7 (Introduction to Networks; Switching, Routing, and Wireless Essentials; Enterprise Networking, Security, and Automation)}{Cisco}{}{}{\vspace{1pt}
Data Analysis Using Python; Python for Data Science.
}}

\end{itemize}

% ============================================================
%     LANGUAGES
% ============================================================

\section{Idiomas}
\vspace{1pt}
\begin{itemize}
\item {Español (nativo), inglés (elemental).}
\end{itemize}

% ============================================================
%     REFERENCES
% ============================================================

\section{Referencias}
\vspace{1pt}
\begin{itemize}
\item{Disponibles bajo petición.}
\end{itemize}

\end{document}
